\section{Discussion}
\label{sec:discussion}

The analysis links a bittide storage requirement to an initial
frequency mismatch through a specified controller and execution
model. Mismatch energy $\delta\omega^\top L\delta\omega$ captures
how frequency differences are distributed over the graph, while
target placement determines the available lower and upper reserve.
Sampling, gains, normalization and graph resistance enter the same
joint test. These are coupled sufficient conditions; a failed test
does not make a configuration unsafe, and maximum frequency spread
alone does not specify the required reserve.

The matched-gain results isolate the implemented nonlinear law from
changes in capacity, gain calibration or protection policy. They do
not establish superiority over an independently optimized PI.
The Fuzzy advantage has two distinct mathematical components.
The exact nonlinear map explains state-dependent increment shaping;
it does not itself order closed-loop performance. The outward
prefix and contracting tail establish the strict excursion ordering.
At fixed capacity the translation lemma then converts this ordering
into a larger admissible target interval for a prescribed frequency
initialization. Full parameter-cell coverage extends the result to
a narrow, quantified interval along the same mismatch direction.
This is a proved advantage for the stated family of configurations,
with positive physical frequencies and no limiter intervention.
It is not a claim that the coarser Lyapunov inner domain is larger
for Fuzzy or that every 400-ppm initialization benefits.

The qualitative experiments directly check these deductions. Their
cases are intentionally chosen around theoretical boundaries and
are not independent statistical discovery. The reference uses the
same continuous held flow as the proof. Callisto separately retains
local scheduling, delay and integer events, so its normal-versus-guard
separation supports the mechanism without transferring the reference
boundary. The upper-target runs identify why receiving occupancy
alone is insufficient for delayed readiness: used credit also includes
forward-transit frames and consumption awaiting reverse feedback,
as expressed by~\eqref{eq:fpi_credit_decomposition}. A physical
buffer below capacity can therefore coexist with a necessary guard.
Fewer guard episodes can accompany lower final data progress.

The occupancy-translation family changes initial occupancy, target
and compatible marking together. It does not certify arbitrary
initial buffer error at a fixed half-capacity target, or preserve
the original application's marking across configurations. Within
each comparison the two controllers have the same marking and
logical configuration. Keeping capacity fixed also preserves
percentage normalization; changing capacity or topology requires
consistent remapping and a fresh certificate.

The nominal reference proves target convergence and syntony with
positive physical progress. Persistent delivery disturbances instead
require finite-window performance measures. Earlier independently
tuned experiments, retained in Appendix~\ref{app:evaluation}, show
benefits in two random regimes and costs in other regimes. They
answer a separate tuning/performance question and do not enlarge
the proved safe domain. Application-level logical synchrony also
requires the prescribed frame relationships; buffer compliance and
frequency agreement do not replace a payload refinement proof.
Extending the invariant analysis to asynchronous sampling, nonzero
link/control delay, integer counts and delayed credit is required
before claiming all-time safety for the full Callisto implementation.

